\documentclass[11pt,a4paper]{article}
\usepackage[utf8]{inputenc}
\usepackage[T1]{fontenc}
\usepackage{amsmath}
\usepackage{amssymb}
\usepackage{geometry}
\usepackage{graphicx}
\usepackage{fancyhdr}
\usepackage{enumitem}
\usepackage{xcolor}
\usepackage{tcolorbox}
\usepackage{hyperref}

% Geometry settings
\geometry{top=1in, bottom=1in, left=1in, right=1in}

% Header and Footer
\pagestyle{fancy}
\fancyhf{}
\rhead{Tashfeen Omran}
\lhead{Homework 15.1: Double Integrals over Rectangles}
\cfoot{\thepage}

% Title
\title{\textbf{Homework 15.1: Double Integrals over Rectangles}}
\author{Tashfeen Omran}
\date{March 2026}

\begin{document}

\maketitle
\tableofcontents
\newpage

\section{Part 1: Comprehensive Introduction, Context, and Prerequisites}

\subsection{Core Concepts}
The topic of \textbf{Double Integrals over Rectangles} extends the concept of the definite integral from single-variable calculus to functions of two variables.
\begin{itemize}
    \item \textbf{The Geometric Interpretation:} Just as the single integral $\int_a^b f(x) \, dx$ represents the \textit{area} under the curve $y=f(x)$, the double integral $\iint_R f(x,y) \, dA$ represents the \textbf{volume} of the solid that lies under the surface $z = f(x,y)$ and above the rectangle $R$ in the $xy$-plane.
    \item \textbf{The Region of Integration ($R$):} In this chapter, we restrict our domain to a rectangle $R = [a,b] \times [c,d] = \{(x,y) \mid a \le x \le b, \, c \le y \le d\}$.
    \item \textbf{Volume as a Limit:} We approximate the volume by dividing the rectangle $R$ into a grid of smaller sub-rectangles, calculating the volume of rectangular prisms (boxes) standing on these sub-rectangles, and summing them up. This is a \textbf{Riemann Sum}.
\end{itemize}

\subsection{Intuition and Derivation}
Imagine a terrain map where $z=f(x,y)$ gives the elevation at coordinate $(x,y)$. To find the volume of the mountain over a rectangular plot of land:
1. We divide the land into a grid of $m \times n$ small squares, each with area $\Delta A = \Delta x \Delta y$.
2. In each small square (sub-rectangle $R_{ij}$), we pick a sample point $(x_{ij}^*, y_{ij}^*)$ and assume the height of the mountain is constant at that point.
3. The volume of the column above that small square is approximately $\text{Height} \times \text{Base Area} = f(x_{ij}^*, y_{ij}^*) \Delta A$.
4. We sum the volumes of all such columns: $V \approx \sum_{i=1}^m \sum_{j=1}^n f(x_{ij}^*, y_{ij}^*) \Delta A$.
5. Taking the limit as the number of squares goes to infinity ($m, n \to \infty$), we get the exact volume:
\[ \iint_R f(x,y) \, dA = \lim_{m,n \to \infty} \sum_{i=1}^m \sum_{j=1}^n f(x_{ij}^*, y_{ij}^*) \Delta A \]

\subsection{Historical Context and Motivation}
In the 18th century, mathematicians like Euler and Lagrange sought to extend the powerful tools of calculus to physical problems in 3D space. Problems such as finding the center of mass of a plate with variable density, the gravitational force exerted by a planet, or the total heat in a metal sheet required integrating over an area rather than a line. This necessitated the development of multiple integration.

\subsection{Key Formulas}
\begin{enumerate}
    \item \textbf{Fubini's Theorem:} This is the practical tool for evaluation. It states that if $f$ is continuous on $R = [a,b] \times [c,d]$, then:
    \[ \iint_R f(x,y) \, dA = \int_a^b \left( \int_c^d f(x,y) \, dy \right) dx = \int_c^d \left( \int_a^b f(x,y) \, dx \right) dy \]
    This allows us to compute a double integral as two sequential single integrals (iterated integrals).
    
    \item \textbf{Separable Functions:} If $f(x,y) = g(x)h(y)$, then:
    \[ \iint_R f(x,y) \, dA = \left( \int_a^b g(x) \, dx \right) \left( \int_c^d h(y) \, dy \right) \]
    
    \item \textbf{Midpoint Rule Approximation:}
    \[ \iint_R f(x,y) \, dA \approx \sum_{i=1}^m \sum_{j=1}^n f(\bar{x}_i, \bar{y}_j) \Delta A \]
    where $\bar{x}_i$ is the midpoint of the interval $[x_{i-1}, x_i]$ and $\bar{y}_j$ is the midpoint of $[y_{j-1}, y_j]$.
\end{enumerate}

\subsection{Prerequisites}
\begin{itemize}
    \item \textbf{Partial Differentiation:} Understanding how to treat $y$ as a constant while working with $x$, and vice versa.
    \item \textbf{Integration Techniques:} $u$-substitution, integration by parts, and power rules are used extensively.
    \item \textbf{Riemann Sums:} Familiarity with the concept of approximating areas with rectangles from Calculus I.
\end{itemize}

\newpage
\section{Part 2: Detailed Homework Solutions}

% PROBLEM 1
\subsection*{Problem 1: Volume of a Solid over a Rectangle}
\textbf{Question:} Evaluate the double integral by first identifying it as the volume of a solid.
\[ \iint_R \sqrt{2} \, dA, \quad R = \{(x, y) \mid 6 \le x \le 8, -1 \le y \le 5\} \]

\textbf{Solution:}
This integral represents the volume of a solid with a constant height $z = f(x,y) = \sqrt{2}$ over the rectangle $R$. This solid is a rectangular box.
The volume of a box is $\text{Area of Base} \times \text{Height}$.

1. Calculate the dimensions of the base rectangle $R$:
   \[ \text{Length (x)} = 8 - 6 = 2 \]
   \[ \text{Width (y)} = 5 - (-1) = 6 \]
2. Calculate the Area of the Base ($A$):
   \[ A = 2 \times 6 = 12 \]
3. Calculate the Volume:
   \[ V = A \times \text{Height} = 12 \times \sqrt{2} = 12\sqrt{2} \]

\textbf{Answer:} $12\sqrt{2}$

% PROBLEM 2
\subsection*{Problem 2: Riemann Sum and Midpoint Rule}
\textbf{Question:} Estimate the volume of the solid that lies below the surface $z = xy$ and above $R = \{(x, y) \mid 6 \le x \le 12, 4 \le y \le 8\}$.
(a) Use a Riemann sum with $m=3, n=2$, sample points at upper right corners.
(b) Use the Midpoint Rule.

\textbf{Solution:}
First, determine the grid sizes:
\[ \Delta x = \frac{12-6}{3} = 2, \quad \Delta y = \frac{8-4}{2} = 2 \]
\[ \Delta A = \Delta x \Delta y = 2 \times 2 = 4 \]

The $x$ sub-intervals are $[6,8], [8,10], [10,12]$.
The $y$ sub-intervals are $[4,6], [6,8]$.

\textbf{(a) Upper Right Corners:}
Sample points $(x_i, y_j)$ are the top-right of each sub-rectangle:
\begin{itemize}
    \item $R_{11}: (8,6) \implies f(8,6) = 48$
    \item $R_{12}: (8,8) \implies f(8,8) = 64$
    \item $R_{21}: (10,6) \implies f(10,6) = 60$
    \item $R_{22}: (10,8) \implies f(10,8) = 80$
    \item $R_{31}: (12,6) \implies f(12,6) = 72$
    \item $R_{32}: (12,8) \implies f(12,8) = 96$
\end{itemize}
Sum of values: $48 + 64 + 60 + 80 + 72 + 96 = 420$.
Volume $\approx 420 \times \Delta A = 420 \times 4 = 1680$.

\textbf{(b) Midpoint Rule:}
Sample points are the centers of the sub-rectangles:
\begin{itemize}
    \item $R_{11}: (7,5) \implies f(7,5) = 35$
    \item $R_{12}: (7,7) \implies f(7,7) = 49$
    \item $R_{21}: (9,5) \implies f(9,5) = 45$
    \item $R_{22}: (9,7) \implies f(9,7) = 63$
    \item $R_{31}: (11,5) \implies f(11,5) = 55$
    \item $R_{32}: (11,7) \implies f(11,7) = 77$
\end{itemize}
Sum of values: $35 + 49 + 45 + 63 + 55 + 77 = 324$.
Volume $\approx 324 \times 4 = 1296$.

\textbf{Answer (a):} $1680$ \\
\textbf{Answer (b):} $1296$

% PROBLEM 3
\subsection*{Problem 3: Estimating Exponential Integral}
\textbf{Question:} Estimate $\iint_R 9xe^{-xy} \, dA$, $R = [0,2] \times [0,1]$, $m=n=2$.
(a) Riemann sum, upper right corners.
(b) Midpoint Rule.

\textbf{Solution:}
$\Delta x = \frac{2-0}{2} = 1$, $\Delta y = \frac{1-0}{2} = 0.5$.
$\Delta A = 1 \times 0.5 = 0.5$.

\textbf{(a) Upper Right Corners:}
Grid points: $x \in \{1, 2\}$, $y \in \{0.5, 1\}$.
\begin{itemize}
    \item $(1, 0.5): f(1, 0.5) = 9(1)e^{-0.5} = 9e^{-0.5} \approx 5.4589$
    \item $(1, 1): f(1, 1) = 9(1)e^{-1} = 9e^{-1} \approx 3.3109$
    \item $(2, 0.5): f(2, 0.5) = 9(2)e^{-1} = 18e^{-1} \approx 6.6218$
    \item $(2, 1): f(2, 1) = 9(2)e^{-2} = 18e^{-2} \approx 2.4360$
\end{itemize}
Sum $\approx 17.8276$.
Integral $\approx 17.8276 \times 0.5 = 8.914$.

\textbf{(b) Midpoint Rule:}
Midpoints: $x \in \{0.5, 1.5\}$, $y \in \{0.25, 0.75\}$.
\begin{itemize}
    \item $(0.5, 0.25): 9(0.5)e^{-0.125} = 4.5 e^{-0.125} \approx 3.9712$
    \item $(0.5, 0.75): 9(0.5)e^{-0.375} = 4.5 e^{-0.375} \approx 3.0928$
    \item $(1.5, 0.25): 9(1.5)e^{-0.375} = 13.5 e^{-0.375} \approx 9.2783$
    \item $(1.5, 0.75): 9(1.5)e^{-1.125} = 13.5 e^{-1.125} \approx 4.3828$
\end{itemize}
Sum $\approx 20.7251$.
Integral $\approx 20.7251 \times 0.5 = 10.363$.

\textbf{Answer (a):} $8.914$ \\
\textbf{Answer (b):} $10.363$

% PROBLEM 4
\subsection*{Problem 4: Iterated Integral Setup and Evaluation}
\textbf{Question:} Surface $z = 8(x^2+y^2)$ over rectangle $R = [0,2] \times [0,3]$.
(a) Set up iterated integral.
(b) Evaluate.

\textbf{Solution:}
\textbf{(a) Setup:}
\[ V = \int_0^3 \int_0^2 8(x^2+y^2) \, dx \, dy \]
(Note: Order $dx \, dy$ or $dy \, dx$ is acceptable).

\textbf{(b) Evaluation:}
Inner integral with respect to $x$:
\[ \int_0^2 8(x^2+y^2) \, dx = 8 \left[ \frac{x^3}{3} + xy^2 \right]_0^2 = 8 \left( \frac{8}{3} + 2y^2 \right) = \frac{64}{3} + 16y^2 \]
Outer integral with respect to $y$:
\[ \int_0^3 \left( \frac{64}{3} + 16y^2 \right) dy = \left[ \frac{64}{3}y + \frac{16y^3}{3} \right]_0^3 \]
\[ = \left( \frac{64}{3}(3) + \frac{16(27)}{3} \right) - 0 = 64 + 16(9) = 64 + 144 = 208 \]

\textbf{Answer:} $208$

% PROBLEM 5
\subsection*{Problem 5: Trigonometric Volume}
\textbf{Question:} $z = 8 \cos x \cos y$ over $R = [-\pi/4, \pi/4] \times [-\pi/4, \pi/4]$.

\textbf{Solution:}
(a) Setup:
\[ \int_{-\pi/4}^{\pi/4} \int_{-\pi/4}^{\pi/4} 8 \cos x \cos y \, dx \, dy \]

(b) Evaluation (using separability):
\[ = 8 \left( \int_{-\pi/4}^{\pi/4} \cos x \, dx \right) \left( \int_{-\pi/4}^{\pi/4} \cos y \, dy \right) \]
Compute one integral:
\[ \int_{-\pi/4}^{\pi/4} \cos t \, dt = [\sin t]_{-\pi/4}^{\pi/4} = \sin(\frac{\pi}{4}) - \sin(-\frac{\pi}{4}) = \frac{\sqrt{2}}{2} - \left(-\frac{\sqrt{2}}{2}\right) = \sqrt{2} \]
Total Volume:
\[ = 8 (\sqrt{2}) (\sqrt{2}) = 8(2) = 16 \]

\textbf{Answer:} $16$

% PROBLEM 6
\subsection*{Problem 6: Riemann Sum Detailed Steps}
\textbf{Question:} $\iint_R 5xe^{-xy} \, dA$, $R=[0,2] \times [0,1]$, $m=n=2$.
(a) Upper Right Corners. Find area of subrectangle, sample points, and sum.
(b) Midpoint Rule. Find sample points.

\textbf{Solution:}
(a)
Area of each subrectangle $\Delta A = \frac{2-0}{2} \times \frac{1-0}{2} = 1 \times 0.5 = 0.5$.

Sample points (Upper Right):
\begin{itemize}
    \item $R_{11}$ ($x \in [0,1], y \in [0,0.5]$): $(1, 0.5)$
    \item $R_{12}$ ($x \in [0,1], y \in [0.5,1]$): $(1, 1)$
    \item $R_{21}$ ($x \in [1,2], y \in [0,0.5]$): $(2, 0.5)$
    \item $R_{22}$ ($x \in [1,2], y \in [0.5,1]$): $(2, 1)$
\end{itemize}

Evaluations:
$f(1, 0.5) = 5e^{-0.5} \approx 3.0327$
$f(1, 1) = 5e^{-1} \approx 1.8394$
$f(2, 0.5) = 10e^{-1} \approx 3.6788$
$f(2, 1) = 10e^{-2} \approx 1.3534$
Sum $\approx 9.9043$.
Riemann Sum $= 9.9043 \times 0.5 = 4.952$.

(b) Sample points (Midpoints):
\begin{itemize}
    \item $R_{11}: (0.5, 0.25)$
    \item $R_{12}: (0.5, 0.75)$
    \item $R_{21}: (1.5, 0.25)$
    \item $R_{22}: (1.5, 0.75)$
\end{itemize}

\textbf{Answer (a):} $4.952$

% PROBLEM 7
\subsection*{Problem 7: Riemann Sum Lower Left}
\textbf{Question:} Estimate volume under $z = 4+x^2+3y$, $R=[1,2] \times [0,3]$. $m=n=2$, Lower Left corners.

\textbf{Solution:}
$\Delta x = 0.5, \Delta y = 1.5$. $\Delta A = 0.75$.
Grid $x$: $[1, 1.5], [1.5, 2]$. Grid $y$: $[0, 1.5], [1.5, 3]$.
Lower Left Points:
\begin{itemize}
    \item $(1, 0): z = 4+1+0 = 5$
    \item $(1, 1.5): z = 4+1+4.5 = 9.5$
    \item $(1.5, 0): z = 4+2.25+0 = 6.25$
    \item $(1.5, 1.5): z = 4+2.25+4.5 = 10.75$
\end{itemize}
Sum $= 5 + 9.5 + 6.25 + 10.75 = 31.5$.
Volume $\approx 31.5 \times 0.75 = 23.625$.

(b) Midpoint Rule:
Points: $(1.25, 0.75), (1.25, 2.25), (1.75, 0.75), (1.75, 2.25)$.
Values:
$f(1.25, 0.75) = 4 + 1.5625 + 2.25 = 7.8125$
$f(1.25, 2.25) = 4 + 1.5625 + 6.75 = 12.3125$
$f(1.75, 0.75) = 4 + 3.0625 + 2.25 = 9.3125$
$f(1.75, 2.25) = 4 + 3.0625 + 6.75 = 13.8125$
Sum $= 43.25$.
Volume $\approx 43.25 \times 0.75 = 32.4375 \approx 32.438$.

\textbf{Answer (a):} $23.625$ \\
\textbf{Answer (b):} $32.438$

% PROBLEM 8
\subsection*{Problem 8: Evaluate Volume}
\textbf{Question:} $\iint_R (4-2y) \, dA$, $R=[0,1] \times [0,1]$.

\textbf{Solution:}
\[ \int_0^1 \int_0^1 (4-2y) \, dy \, dx \]
Inner (wrt $y$): $[4y - y^2]_0^1 = (4-1) = 3$.
Outer (wrt $x$): $\int_0^1 3 \, dx = [3x]_0^1 = 3$.

\textbf{Answer:} $3$

% PROBLEM 9
\subsection*{Problem 9: Partial Integrals}
\textbf{Question:} $f(x,y) = 5x + 3x^2y^2$. Find $\int_0^2 f(x,y) \, dx$ and $\int_0^3 f(x,y) \, dy$.

\textbf{Solution:}
1. $\int_0^2 (5x + 3x^2y^2) \, dx = [\frac{5x^2}{2} + x^3y^2]_0^2 = \frac{5(4)}{2} + (8)y^2 = 10 + 8y^2$.
2. $\int_0^3 (5x + 3x^2y^2) \, dy = [5xy + 3x^2\frac{y^3}{3}]_0^3 = [5xy + x^2y^3]_0^3 = 5x(3) + x^2(27) = 15x + 27x^2$.

\textbf{Answer:} $\int_0^2 f dx = 10+8y^2$, $\int_0^3 f dy = 15x+27x^2$.

% PROBLEM 10
\subsection*{Problem 10: Iterated Integral}
\textbf{Question:} $\int_1^2 \int_0^3 (6x^2y - 4x) \, dy \, dx$.

\textbf{Solution:}
Inner ($dy$):
\[ \int_0^3 (6x^2y - 4x) \, dy = \left[ 3x^2y^2 - 4xy \right]_0^3 \]
\[ = (3x^2(9) - 4x(3)) - 0 = 27x^2 - 12x \]
Outer ($dx$):
\[ \int_1^2 (27x^2 - 12x) \, dx = \left[ 9x^3 - 6x^2 \right]_1^2 \]
\[ = (9(8) - 6(4)) - (9(1) - 6(1)) \]
\[ = (72 - 24) - (3) = 48 - 3 = 45 \]

\textbf{Answer:} $45$

% PROBLEM 11
\subsection*{Problem 11: Iterated Integral}
\textbf{Question:} $\int_0^4 \int_0^1 (x+y)^2 \, dx \, dy$.

\textbf{Solution:}
Inner ($dx$):
\[ \int_0^1 (x+y)^2 \, dx = \left[ \frac{(x+y)^3}{3} \right]_{x=0}^{x=1} = \frac{(1+y)^3}{3} - \frac{y^3}{3} \]
Outer ($dy$):
\[ \int_0^4 \left( \frac{1}{3}(1+y)^3 - \frac{1}{3}y^3 \right) dy = \frac{1}{3} \left[ \frac{(1+y)^4}{4} - \frac{y^4}{4} \right]_0^4 \]
\[ = \frac{1}{12} \left[ ((1+4)^4 - 4^4) - (1^4 - 0^4) \right] \]
\[ = \frac{1}{12} [ (625 - 256) - 1 ] = \frac{1}{12} [ 369 - 1 ] = \frac{368}{12} = \frac{92}{3} \]

\textbf{Answer:} $92/3$

% PROBLEM 12
\subsection*{Problem 12: Order dy dx}
\textbf{Question:} $\int_{-6}^1 \int_1^2 (x^2 + y^{-2}) \, dy \, dx$.

\textbf{Solution:}
Inner ($dy$):
\[ \int_1^2 (x^2 + y^{-2}) \, dy = \left[ x^2y - y^{-1} \right]_1^2 = (2x^2 - \frac{1}{2}) - (x^2 - 1) = x^2 + \frac{1}{2} \]
Outer ($dx$):
\[ \int_{-6}^1 (x^2 + \frac{1}{2}) \, dx = \left[ \frac{x^3}{3} + \frac{x}{2} \right]_{-6}^1 \]
\[ = (\frac{1}{3} + \frac{1}{2}) - (\frac{-216}{3} + \frac{-6}{2}) = \frac{5}{6} - (-72 - 3) = \frac{5}{6} + 75 = \frac{455}{6} \]

\textbf{Answer:} $455/6$

% PROBLEM 13
\subsection*{Problem 13: Trig Integral}
\textbf{Question:} $\int_{-6}^6 \int_0^{\pi/2} (y + y^2 \cos x) \, dx \, dy$.

\textbf{Solution:}
Inner ($dx$):
\[ \int_0^{\pi/2} (y + y^2 \cos x) \, dx = [yx + y^2 \sin x]_0^{\pi/2} \]
\[ = (y(\frac{\pi}{2}) + y^2(1)) - 0 = \frac{\pi}{2}y + y^2 \]
Outer ($dy$):
\[ \int_{-6}^6 (\frac{\pi}{2}y + y^2) \, dy \]
Note that $\frac{\pi}{2}y$ is an odd function on a symmetric interval $[-6,6]$, so its integral is 0.
\[ \int_{-6}^6 y^2 \, dy = \left[ \frac{y^3}{3} \right]_{-6}^6 = \frac{216}{3} - \frac{-216}{3} = 72 + 72 = 144 \]

\textbf{Answer:} $144$

% PROBLEM 14
\subsection*{Problem 14: Logarithmic Integral}
\textbf{Question:} $\int_1^9 \int_1^8 \frac{\ln y}{xy} \, dy \, dx$.

\textbf{Solution:}
Separate variables: $\left( \int_1^9 \frac{1}{x} \, dx \right) \left( \int_1^8 \frac{\ln y}{y} \, dy \right)$.
1. $\int_1^9 \frac{1}{x} \, dx = [\ln |x|]_1^9 = \ln 9$.
2. $\int_1^8 \frac{\ln y}{y} \, dy$. Let $u = \ln y, du = \frac{1}{y}dy$. Bounds $u(1)=0, u(8)=\ln 8$.
   $\int_0^{\ln 8} u \, du = [\frac{u^2}{2}]_0^{\ln 8} = \frac{(\ln 8)^2}{2}$.
Total: $\frac{\ln 9 (\ln 8)^2}{2}$.
Note: $\ln 9 = 2 \ln 3$, $\ln 8 = 3 \ln 2$.
Result $= \frac{2 \ln 3 (3 \ln 2)^2}{2} = \ln 3 (9 (\ln 2)^2) = 9 \ln 3 (\ln 2)^2$.

\textbf{Answer:} $\frac{(\ln 9)(\ln 8)^2}{2}$ or $9 \ln 3 (\ln 2)^2$

% PROBLEM 15
\subsection*{Problem 15: Separable Trig}
\textbf{Question:} $\int_0^9 \int_0^{\pi/2} t^2 \sin^3 \phi \, d\phi \, dt$.

\textbf{Solution:}
Separate: $(\int_0^9 t^2 dt)(\int_0^{\pi/2} \sin^3 \phi d\phi)$.
1. $\int_0^9 t^2 dt = [\frac{t^3}{3}]_0^9 = \frac{729}{3} = 243$.
2. $\int \sin^3 \phi \, d\phi = \int (1-\cos^2 \phi)\sin \phi \, d\phi$. Let $u=\cos \phi, du=-\sin \phi d\phi$.
   $\int_1^0 (1-u^2)(-du) = \int_0^1 (1-u^2) du = [u - \frac{u^3}{3}]_0^1 = 1 - 1/3 = 2/3$.
Total: $243 \times \frac{2}{3} = 162$.

\textbf{Answer:} $162$

% PROBLEM 16
\subsection*{Problem 16: Algebraic Integral}
\textbf{Question:} $\int_0^1 \int_0^1 \frac{7xy}{\sqrt{x^2+y^2}} \, dy \, dx$.

\textbf{Solution:}
Inner ($dy$): Let $u = x^2+y^2, du = 2y dy \implies \frac{1}{2}du = y dy$.
$\int \frac{7x}{2\sqrt{u}} du = 7x \sqrt{u} = 7x\sqrt{x^2+y^2}$.
Evaluate at $y=0, y=1$:
$[7x\sqrt{x^2+1} - 7x\sqrt{x^2}] = 7x\sqrt{x^2+1} - 7x^2$.
Outer ($dx$):
$\int_0^1 (7x(x^2+1)^{1/2} - 7x^2) dx$.
Part 1: $\int 7x(x^2+1)^{1/2} dx$. Let $v=x^2+1, dv=2x dx$. $\frac{7}{2} \int v^{1/2} dv = \frac{7}{2} \cdot \frac{2}{3} v^{3/2} = \frac{7}{3}(x^2+1)^{3/2}$.
Part 2: $\int 7x^2 dx = \frac{7x^3}{3}$.
Evaluate from 0 to 1:
$[ \frac{7}{3}(2)^{3/2} - \frac{7}{3}(1) ] - [ \frac{7}{3}(1)^{3/2} - 0 ]$.
$= \frac{7}{3}(2\sqrt{2}) - \frac{7}{3} - \frac{7}{3} = \frac{14\sqrt{2} - 14}{3}$.

\textbf{Answer:} $\frac{14(\sqrt{2}-1)}{3}$

% PROBLEM 17
\subsection*{Problem 17: Secant Integral}
\textbf{Question:} $\iint_R x \sec^2 y \, dA$, $R = [0,8] \times [0, \pi/4]$.

\textbf{Solution:}
Separate: $(\int_0^8 x dx)(\int_0^{\pi/4} \sec^2 y dy)$.
1. $[\frac{x^2}{2}]_0^8 = 32$.
2. $[\tan y]_0^{\pi/4} = 1 - 0 = 1$.
Total: $32 \times 1 = 32$.

\textbf{Answer:} $32$

% PROBLEM 18
\subsection*{Problem 18: Double Integral}
\textbf{Question:} $\iint_R (6y + xy^{-2}) \, dA$, $R = [0,2] \times [1,2]$.

\textbf{Solution:}
Inner ($dx$): $[6xy + \frac{x^2}{2}y^{-2}]_0^2 = 12y + 2y^{-2}$.
Outer ($dy$): $\int_1^2 (12y + 2y^{-2}) dy = [6y^2 - 2y^{-1}]_1^2$.
$= (6(4) - 1) - (6(1) - 2) = (24-1) - 4 = 23 - 4 = 19$.

\textbf{Answer:} $19$

% PROBLEM 19
\subsection*{Problem 19: Rational Function}
\textbf{Question:} $\iint_R \frac{3xy^2}{x^2+1} \, dA$, $R = [0,1] \times [-2,2]$.

\textbf{Solution:}
Separate: $(\int_0^1 \frac{3x}{x^2+1} dx)(\int_{-2}^2 y^2 dy)$.
1. Let $u=x^2+1, du=2x dx$. $\frac{3}{2} \int \frac{1}{u} du = \frac{3}{2} \ln|u|_1^2 = \frac{3}{2} \ln 2$.
2. $[\frac{y^3}{3}]_{-2}^2 = \frac{8}{3} - \frac{-8}{3} = \frac{16}{3}$.
Total: $\frac{3}{2} \ln 2 \times \frac{16}{3} = 8 \ln 2$.

\textbf{Answer:} $8 \ln 2$

% PROBLEM 20
\subsection*{Problem 20: Similar Rational Function}
\textbf{Question:} $\iint_R \frac{4xy^2}{x^2+1} \, dA$, $R = [0,4] \times [-3,3]$.

\textbf{Solution:}
Separate: $(\int_0^4 \frac{4x}{x^2+1} dx)(\int_{-3}^3 y^2 dy)$.
1. $2 \ln(x^2+1)|_0^4 = 2 \ln 17$.
2. $[\frac{y^3}{3}]_{-3}^3 = 9 - (-9) = 18$.
Total: $36 \ln 17$.

\textbf{Answer:} $36 \ln 17$

% PROBLEM 21
\subsection*{Problem 21: Integration by Parts}
\textbf{Question:} $\iint_R 11x \sin(x+y) \, dA$, $R = [0, \pi/6] \times [0, \pi/3]$.

\textbf{Solution:}
Inner ($dy$):
$\int_0^{\pi/3} 11x \sin(x+y) dy = [-11x \cos(x+y)]_0^{\pi/3} = -11x [\cos(x+\pi/3) - \cos x]$.
Outer ($dx$):
$11 \int_0^{\pi/6} x (\cos x - \cos(x+\pi/3)) dx$.
Use $\int x \cos(x+C) dx = x \sin(x+C) + \cos(x+C)$.
Term 1: $\int x \cos x dx = [x \sin x + \cos x]_0^{\pi/6} = (\frac{\pi}{6}\frac{1}{2} + \frac{\sqrt{3}}{2}) - 1$.
Term 2: $\int x \cos(x+\pi/3) dx = [x \sin(x+\pi/3) + \cos(x+\pi/3)]_0^{\pi/6}$.
Upper: $\frac{\pi}{6} \sin(\pi/2) + \cos(\pi/2) = \frac{\pi}{6} + 0$.
Lower: $0 + \cos(\pi/3) = 1/2$.
Result Term 2: $\frac{\pi}{6} - 1/2$.
Combine: $11 [ (\frac{\pi}{12} + \frac{\sqrt{3}}{2} - 1) - (\frac{\pi}{6} - \frac{1}{2}) ]$.
$= 11 [ \frac{\pi}{12} - \frac{2\pi}{12} + \frac{\sqrt{3}}{2} - 1 + \frac{1}{2} ] = 11 [ -\frac{\pi}{12} + \frac{\sqrt{3}}{2} - \frac{1}{2} ]$.

\textbf{Answer:} $11(\frac{\sqrt{3}-1}{2} - \frac{\pi}{12})$

% PROBLEM 22
\subsection*{Problem 22: Log Integration}
\textbf{Question:} $\iint_R \frac{2x}{1+xy} \, dA$, $R = [0,2] \times [0,1]$.

\textbf{Solution:}
Inner ($dy$): $\int_0^1 \frac{2x}{1+xy} dy$. Let $u=1+xy$, $du=x dy$.
$= 2 \int \frac{1}{u} du = 2 [\ln(1+xy)]_0^1 = 2 \ln(1+x)$.
Outer ($dx$): $\int_0^2 2 \ln(1+x) dx$.
Let $w=1+x$. $\int_1^3 2 \ln w dw$.
IBP: $\int \ln w = w \ln w - w$.
$2 [ w \ln w - w ]_1^3 = 2 [ (3 \ln 3 - 3) - (-1) ] = 2 (3 \ln 3 - 2) = 6 \ln 3 - 4$.

\textbf{Answer:} $6 \ln 3 - 4$

\newpage
\section{Part 3: In-Depth Analysis of Problems and Techniques}

\subsection{A) Problem Types and General Approach}
\begin{enumerate}
    \item \textbf{Geometric Volume Problems (Q1, Q8):} These require identifying the integral as the volume of a simple solid (like a box or prism). Strategy: Calculate $Area(Base) \times Height$ or use basic geometry formulas instead of integration.
    \item \textbf{Approximation Problems (Q2, Q3, Q6, Q7):} These ask for Riemann Sums or Midpoint Rules. Strategy: Accurately calculate $\Delta x, \Delta y$, determine grid points, and sum $f(x_i, y_j) \Delta A$.
    \item \textbf{Separable Integrals (Q5, Q14, Q15, Q17, Q19, Q20):} The integrand is $f(x)g(y)$. Strategy: Split into $(\int f(x)dx)(\int g(y)dy)$ to save time.
    \item \textbf{Iterated Integrals (Q4, Q10, Q11, Q12, Q13, Q16, Q18, Q21, Q22):} These require sequential integration. Strategy: Integrate the inner variable first (treating outer as constant), substitute bounds, then integrate the resulting single-variable function.
\end{enumerate}

\subsection{B) Key Algebraic and Calculus Manipulations}
\begin{itemize}
    \item \textbf{Separation of Variables:} Used in Q5 ($8 \cos x \cos y$) to turn a double integral into a product of two easy single integrals.
    \item \textbf{Power Rule for Negative Exponents:} Used in Q12 ($y^{-2} \to -y^{-1}$) and Q18 ($xy^{-2}$). Crucial for ensuring the correct sign during evaluation.
    \item \textbf{Integration by Parts:} Required in Q21 ($x \sin(x+y)$) and Q22 ($\ln(1+x)$). This handles products of algebraic and transcendental functions.
    \item \textbf{Partial Fractions / Logarithmic Integration:} Used in Q19/Q20 ($ \int \frac{x}{x^2+1} dx = \frac{1}{2}\ln(x^2+1) $). A common pattern where the numerator is the derivative of the denominator.
\end{itemize}

\section{Part 4: ``Cheatsheet'' and Tips for Success}

\begin{tcolorbox}[title=Essential Formulas]
    \textbf{Fubini's Theorem:} $\iint_R f(x,y) dA = \int_a^b \int_c^d f(x,y) dy dx$ \\
    \textbf{Separability:} If $f(x,y) = g(x)h(y)$, then $\iint f dA = (\int g dx)(\int h dy)$ \\
    \textbf{Midpoint Rule:} Use midpoints $\bar{x}, \bar{y}$ of intervals. Error is usually lower than corner approximations.
\end{tcolorbox}

\begin{itemize}
    \item \textbf{Shortcut:} Check if the function is separable immediately. If yes, you avoid managing nested variables.
    \item \textbf{Trap:} Forgetting to treat the "other" variable as a constant. In $\int x^2 y dy$, $x^2$ is just a number like 5.
    \item \textbf{Trap:} Order of Integration. If you write $dy dx$, the inner bounds must correspond to $y$ and outer to $x$.
    \item \textbf{Tip:} Symmetry. If integrating an odd function (like $\sin x$ or $x^3$) over a symmetric interval $[-a, a]$, the result is 0 (seen in Q13).
\end{itemize}

\section{Part 5: Conceptual Synthesis and The ``Big Picture''}
\subsection{Thematic Connections}
The central theme is \textbf{accumulation}. Just as a single integral accumulates height to find area, a double integral accumulates area slices to find volume. This concept generalizes to Triple Integrals (accumulating volume to find mass/hyper-volume) later in the course.

\subsection{Forward and Backward Links}
\textbf{Backward:} Relies heavily on Calc I/II integration techniques (u-sub, parts).
\textbf{Forward:} This is the foundation for Surface Integrals and Flux in Vector Calculus. Understanding how to parameterize a domain $R$ is crucial for Green's Theorem and Stokes' Theorem.

\section{Part 6: Real-World Application and Modeling}
\subsection{A) Concrete Scenarios}
1. \textbf{Civil Engineering (Earthworks):} Determining the volume of soil to be removed from a plot of land for a foundation. The terrain height is $z=f(x,y)$.
2. \textbf{Meteorology:} Calculating the average rainfall over a rectangular state/region. Average value formula: $\frac{1}{Area} \iint f(x,y) dA$.
3. \textbf{Finance (Joint Probability):} Calculating the probability that two correlated stock prices $X$ and $Y$ fall within certain ranges. This requires integrating the Joint Probability Density Function (PDF) $f(x,y)$ over a rectangle.

\subsection{B) Model Problem Setup (Finance)}
\textbf{Scenario:} An investor wants to know the probability that Stock A is between \$100-\$110 and Stock B is between \$50-\$60.
\textbf{Model:} Let $x$ be price of A, $y$ be price of B.
The Joint PDF is modeled as $f(x,y) = k e^{-(x+y)/100}$ for $x,y > 0$.
\textbf{Integral:}
\[ P(100 \le x \le 110, 50 \le y \le 60) = \int_{100}^{110} \int_{50}^{60} k e^{-(x+y)/100} \, dy \, dx \]

\section{Part 7: Common Variations and Untested Concepts}
\textbf{Concept:} \textbf{Average Value of a Function.}
Your homework calculated volumes, but often we want the average height.
Formula: $f_{ave} = \frac{1}{\text{Area}(R)} \iint_R f(x,y) dA$.
\textbf{Example:} Find average of $f(x,y)=xy$ over $R=[0,2]\times[0,2]$.
$\iint xy dA = \int_0^2 x dx \int_0^2 y dy = 2 \times 2 = 4$.
Area $= 2 \times 2 = 4$.
Average $= 4/4 = 1$.

\section{Part 8: Advanced Diagnostic Testing: ``Find the Flaw''}

\subsection*{Flawed Problem 1: Variable Handling}
\textbf{Problem:} Evaluate $\int_0^1 \int_0^1 x^2 y \, dx \, dy$.
\textbf{Flawed Solution:}
$\int_0^1 x^2 y \, dx = [\frac{x^3}{3} \frac{y^2}{2}]_0^1 = \frac{1}{6}$.
$\int_0^1 \frac{1}{6} dy = [\frac{1}{6}y]_0^1 = 1/6$.
\textbf{Error:} The student integrated both $x$ and $y$ simultaneously in the inner step.
\textbf{Correction:} In the inner integral $\int x^2 y dx$, $y$ is constant. Correct step: $[ \frac{x^3}{3}y ]_0^1 = \frac{1}{3}y$. Then integrate $\frac{1}{3}y$ wrt $y$. Result: $1/6$ (Answer accidentally correct, method wrong).

\subsection*{Flawed Problem 2: Separability Abuse}
\textbf{Problem:} Evaluate $\iint (x+y) \, dA$.
\textbf{Flawed Solution:}
$= (\int x \, dx)(\int y \, dy)$.
\textbf{Error:} $x+y$ is not a product $g(x)h(y)$, so it cannot be separated into a product of integrals.
\textbf{Correction:} Integrate terms linearly: $\iint x dA + \iint y dA$.

\subsection*{Flawed Problem 3: Trig Evaluation}
\textbf{Problem:} $\int_0^{\pi} \cos x \, dx$.
\textbf{Flawed Solution:} $[\sin x]_0^{\pi} = \sin \pi - \sin 0 = 1 - 0 = 1$.
\textbf{Error:} $\sin \pi = 0$, not $1$.
\textbf{Correction:} Result is $0$.

\subsection*{Flawed Problem 4: Integration by Parts Sign}
\textbf{Problem:} $\int x \cos x \, dx$.
\textbf{Flawed Solution:} $x \sin x - \int \sin x \, dx = x \sin x - \cos x$.
\textbf{Error:} The integral of $\sin x$ is $-\cos x$. The negative sign subtracts a negative.
\textbf{Correction:} $x \sin x - (-\cos x) = x \sin x + \cos x$.

\subsection*{Flawed Problem 5: Logarithm Bounds}
\textbf{Problem:} $\int_{-2}^{-1} \frac{1}{x} dx$.
\textbf{Flawed Solution:} $[\ln x]_{-2}^{-1} = \ln(-1) - \ln(-2)$. Undefined.
\textbf{Error:} The integral of $1/x$ is $\ln|x|$.
\textbf{Correction:} $\ln|-1| - \ln|-2| = 0 - \ln 2 = -\ln 2$.

\end{document}